%% kapa: exam=1 round=19 year=2008 period=1 form=A subject=civil_law subject_ko=민법 seq=36 number=36
%% source: Q-Net 감정평가사 기출문제 (cst00309) artlSeq=1001962 file=extracted/1008/1001962/2008년 제19회 감정평가사 1차시험 기출문제/1교시.hwp
%% conversion: kordoc 4.24.51 -> tools/gen1.py (auto); review=proofread-sonnet; flags=none
\begin{question}{36}
전세권에 관한 설명으로 틀린 것은? (다툼이 있는 경우에는 판례에 의함)
\begin{choices}
  \item 건물의 일부에 대하여만 전세권이 설정된 경우라도 전세권자는 건물 전부에 대하여 후순위권리자보다 전세금의 우선변제를 받을 권리가 있다.
  \item 전세권의 존속 중이라도 장래 그 전세권이 소멸하는 경우에 전세금반환청구권이 발생하는 것을 조건으로 한 그 장래의 조건부채권을 양도할 수 있다.
  \item 타인의 토지에 있는 건물에 전세권을 설정해 준 자는 전세권자의 동의 없이 그 건물의 소유를 목적으로 하는 자신의 지상권을 소멸시키지 못한다.
  \item 전세금액은 필요적 등기사항으로서, 등기된 금액에 대하여는 제3자에게 대항할 수 있다.
  \item 1번 저당권이 설정된 후 전세권이 설정되고 그 후 2번 저당권이 설정된 경우, 2번 저당권자가 경매를 신청하여 전세목적물이 매각되었다면 전세권은 소멸하지 않는다.
\end{choices}
\end{question}
\answer{5}
